\chapter{Packets, types, holonomy, and restricted realization}\label{ch:packets}
\index{spherical packet}
\index{type packet}
\index{holonomy packet}
\index{realization fiber}
\index{quantum Latin square}
\index{Langlands-style packet}


\section{Why coarse spectra are not enough}
A finite symmetry problem often comes with a natural commuting operator algebra.  Its joint spectrum is attractive because it is canonical, compact, and computationally cheap.  Yet the joint eigenvalues and their multiplicities can be identical across objects whose realization behavior is different.

This chapter develops the hierarchy
\begin{equation}\label{eq:packet-hierarchy}
\text{spherical packet}
\prec
\text{embedded type packet}
\prec
\text{holonomy packet}
\prec
\text{realization fiber}.
\end{equation}
The relation $\prec$ means ``retains strictly more potentially relevant information,'' not that every refinement is complete.

\begin{figure}[ht]
\centering
\includegraphics[width=.95\textwidth]{packet_hierarchy.pdf}
\caption{The packet hierarchy used to organize finite threshold systems.}
\label{fig:packet-realization-hierarchy}
\end{figure}

\section{Spherical packets}
Let $T_1,\ldots,T_m$ be commuting normal operators on a finite-dimensional Hilbert space.  Their unique coarsest joint spectral resolution is
\[
T_j=\sum_{\lambda}t_j(\lambda)B_\lambda,
\qquad
\sum_\lambda B_\lambda=I.
\]
The tuple
\[
\lambda=(t_1(\lambda),\ldots,t_m(\lambda))
\]
is the spherical label.  The abstract spectral packet records these labels and the ranks of $B_\lambda$.

Two different embedded systems can have the same abstract packet.  The actual matrices $B_\lambda$, endpoint coefficients, and transition operators may still differ.

\begin{principle}[Spherical collapse]
Passing from embedded operators to eigenvalues and multiplicities can erase exactly the incidence or transition data that controls restricted realization.
\end{principle}

\section{The Latin-square association algebra}
A Latin square $L$ of order $n$ has $n^2$ cells.  Let $R,C,S$ be adjacency matrices for equality of row, column, and symbol.  The algebra generated by $I,J,R,C,S$ is commutative and its abstract multiplication table depends only on $n$, not on the paratopy class of $L$.

Define
\[
Q_R=\frac{I+R}{n},
\qquad
Q_C=\frac{I+C}{n},
\qquad
Q_S=\frac{I+S}{n},
\qquad
B_0=\frac{J}{n^2}.
\]
Then the five joint projectors are
\[
B_0,
\quad Q_R-B_0,
\quad Q_C-B_0,
\quad Q_S-B_0,
\quad I-Q_R-Q_C-Q_S+2B_0.
\]
Their multiplicities are
\[
1,
\quad n-1,
\quad n-1,
\quad n-1,
\quad (n-1)(n-2).
\]
The Latin-square graph $R+C+S$ has universal spectrum
\[
3(n-1)^1,
\qquad
(n-3)^{3(n-1)},
\qquad
(-3)^{(n-1)(n-2)}.
\]
At $n=6$ this becomes
\[
15^1,\qquad3^{15},\qquad(-3)^{20}.
\]
Thus the graph spectrum cannot distinguish the twelve main classes of order six.

\begin{computedbox}
The accompanying atlas transcribes all twelve order-six main classes used in the quantum-Latin-square reduction and verifies that their abstract $R,C,S$ joint eigenvalue packets coincide.  This is an exact no-go result for an eigenvalue-only explanation of the order-six obstruction.
\end{computedbox}

\section{Transition fusion also collapses}
Let $\rho_i$ be the row permutations of a Latin square.  The symmetrized aggregate of all pair transitions satisfies
\begin{equation}\label{eq:transition-fusion}
\sum_{i<j}
\left(
P_{\rho_i\rho_j^{-1}}
+
P_{\rho_j\rho_i^{-1}}
\right)
=n(J-I).
\end{equation}
Hence first-order fusion destroys the class information as well.

The lesson is general: summing typed transitions can be as destructive as discarding the embedded projectors.  Order and labels matter.

\section{Embedded type data}
A type packet retains the actual projectors, support partitions, distinguished incidence sets, or local representation types.  Examples include:
\begin{itemize}
\item the row, column, and symbol partitions of a Latin square;
\item a projective-line closure system;
\item the flag stabilized by a circuit subgroup;
\item support patterns of unitary or quantum-Latin entries;
\item ramification or block data in a finite group action.
\end{itemize}

In the order-six Latin-square computation, the discrete projective algorithm leaves two hard classes with a spread of twelve disjoint projective lines on the $36$ cells.  The spread adjacency operator commutes with $R,C,S$ and therefore refines the spherical packet.  The two hard classes nevertheless retain the same refined abstract spectrum.  One type operator is not complete.

\section{Holonomy packets}
Let a family of local charts or rows be represented by permutations $\rho_i$.  Define normalized transitions
\[
h_{ij}=\rho_i\rho_j^{-1}
\]
and the transition group
\[
H_r(L)=\langle h_{ij}:i,j\rangle.
\]
A holonomy packet may include:
\begin{itemize}
\item the isomorphism type or order of $H_r(L)$,
\item the cycle-type inventory of the distinguished $h_{ij}$,
\item block systems and orbit partitions,
\item characters or matrix coefficients of words in the $h_{ij}$,
\item noncommutative moments
\[
\operatorname{tr}(h_{i_1j_1}\cdots h_{i_kj_k}),
\]
\item the corresponding packets for row, column, and symbol roles.
\end{itemize}

Isotopy conjugates these transition groups; paratopy permutes the three roles.  Therefore the unordered multiset of role signatures is a paratopy invariant.

\begin{computedbox}
For the twelve order-six main classes, the finite packet consisting of transition-group order, pair-transition cycle inventory, and nontrivial block sizes is invariant under tested isotopies and all six role permutations.  The twelve packets are pairwise distinct.
\end{computedbox}
This classification result does not prove orthonormal representability.  It demonstrates that noncommutative transition data recover information lost under spherical fusion.

\section{Imprimitivity as hidden decomposition}
A transitive permutation group is imprimitive when it preserves a nontrivial block system.  In combinatorial designs, imprimitivity often signals a subsquare, product, congruence, or module decomposition.

For a Latin square of order $2m$, the normalized row-transition group preserves a two-block system of sizes $m+m$ precisely when the square contains an $m\times m$ Latin subsquare, under the normalization used in the atlas.  At order six this identifies the four classes containing $3\times3$ subsquares.

\begin{principle}[Clifford-type decomposition]
Combinatorial substructure can be detected as imprimitivity of a transition packet.  The spherical packet sees aggregate multiplicities; the holonomy packet sees how those multiplicities are assembled into blocks.
\end{principle}

\section{Realization fibers}
Let $P$ denote packet data and let
\[
\operatorname{Real}_{\mathrm{res}}(P)
\]
be the space of restricted realizations compatible with $P$.  The packet may completely classify the discrete object while the realization fiber remains empty, positive dimensional, disconnected, or singular.

Examples of realization conditions include:
\begin{itemize}
\item a flag-preserving geodesic with a declared length;
\item a product-state orthonormal representation;
\item a positive semidefinite Gram matrix with a zero pattern;
\item a tensor factorization;
\item a global descent object compatible with local packets;
\item a circuit satisfying ancilla restoration on every input.
\end{itemize}

Thus classification and realization are distinct layers.

\section{Quantum Latin squares as an existence defect}
A quantum Latin square is an array \cite{MustoVicary2016} of vectors whose rows and columns are orthonormal bases.  Ball and Simoens prove that a pair of non-entangled mutually orthogonal quantum Latin squares of order six does not exist, while earlier work provides entangled order-six solutions and hence AME$(4,6)$ states \cite{BallSimoens2026,RatherEtAl2022,RatherEtAl2023}.

This is naturally expressed as
\[
\operatorname{Real}_{\mathrm{product}}(P_6)=\varnothing,
\qquad
\operatorname{Real}_{\mathrm{entangled}}(P_6)\ne\varnothing.
\]
It is an existence defect, not a gate-count defect.

\begin{figure}[ht]
\centering
\includegraphics[width=.92\textwidth]{quantum_latin_realization_gap.pdf}
\caption{The order-six quantum-Latin distinction is a restricted-versus-ambient realization gap.}
\label{fig:quantum-latin-gap}
\end{figure}

The source proof uses support patterns, a reduction to one classical square, orthonormal representations of complements of Latin-square graphs, projective-line closure, finite case analysis over twelve main classes, and a continuous argument for the two remaining subsquare cases.  No theorem in the unitriangular or Clifford chapters is inferred from this result.

\section{First-arrival packets}
For a finite generator set $S$, the endpoint Green kernel
\[
\mathcal G_{S,g}(z)
=
\langle\delta_g,(I-zA_S)^{-1}\delta_e\rangle
\]
contains more information than the spectrum of $A_S$.  Its first nonzero degree is the distance to $g$.

If
\[
A_S=\sum_\lambda \lambda B_\lambda,
\]
then
\[
\mathcal G_{S,g}(z)
=
\sum_\lambda
\frac{\langle\delta_g,B_\lambda\delta_e\rangle}{1-\lambda z}.
\]
The endpoint matrix coefficients of the embedded projectors determine cancellation between spectral components.  This is the metric analogue of retaining a type rather than only an unramified eigencharacter.

\section{Typed Green kernels}
Let $H\le G$ and let $\mu_k(h)$ count ambient words of length $k$ with holonomy $h\in H$.  For an irreducible representation $\tau$ of $H$, define
\[
\widehat\mu_k(\tau)
=
\sum_{h\in H}\mu_k(h)\tau(h)
\]
and
\[
\widehat{\mathcal G}_\tau(z)
=
\sum_{k\ge0}\widehat\mu_k(\tau)z^k.
\]
Fourier inversion recovers every $\mu_k(h)$.  The trivial representation records only total return to the subgroup; nontrivial types resolve the holonomy.

\begin{provedbox}
For finite $H$, the full family of typed Green kernels over $\widehat H$ determines every endpoint Green series and therefore every ambient word distance from the identity to an element of $H$.
\end{provedbox}

\section{Operational Langlands analogy}
The hierarchy resembles the local organization of representation theory:
\[
\text{spherical character}
\to
\text{type}
\to
\text{ramification or monodromy}
\to
\text{realization and descent}.
\]
The analogy is useful because it prevents coarse spectral data from being mistaken for a complete correspondence.  It does not imply that global Langlands reciprocity predicts either the circuit threshold or the quantum-Latin result.

\begin{boundarybox}
The word ``Langlands'' in this monograph refers to an operational packet architecture: commuting local spectra, type refinement, noncommutative holonomy, tensor assembly, and descent obligations.  No general automorphic or Galois correspondence is claimed or used.
\end{boundarybox}

\section{A general finite threshold system}
A packet-resolved threshold system consists of:
\begin{enumerate}
\item a finite group, groupoid, or association scheme;
\item a commutative spherical algebra with exact projectors;
\item embedded incidence or type data;
\item a transition group or monodromy representation;
\item restricted and ambient realization categories;
\item a metric, multiplicity, or descent obstruction.
\end{enumerate}

The first threshold may then be defined as
\[
\tau=
\min\{n:\exists P\text{ with }c_{\mathrm{res}}(P)>c_{\mathrm{amb}}(P)\}.
\]
The cost can be finite or infinite.

\section{Open problems}
\begin{question}
For which finite association schemes does the full embedded Bose--Mesner algebra determine the restricted realization fiber?
\end{question}
\begin{question}
What bounded degree of noncommutative transition moments separates all Latin squares of a fixed order?
\end{question}
\begin{question}
Can typed Green kernels yield new lower bounds without enumerating the full ambient group?
\end{question}
\begin{question}
When does imprimitivity of a transition packet characterize a tensor or subsquare decomposition in quantum combinatorial designs?
\end{question}
